% In this file you should put all LaTeX macros and settings to be used both by
% the pdf version and the web version.
% This should be most of your macros.

% The theorem-like environments defined below are those that appear by default
% in the dependency graph. See the README of leanblueprint if you need help to
% customize this. 
% The configuration below use the theorem counter for all those environments
% (this is what the [theorem] arguments mean) and never resets it.
% If you want for instance to number them within chapters then you can add
% [chapter] at the end of the next line.
\newtheorem{theorem}{Theorem}
\newtheorem{proposition}[theorem]{Proposition}
\newtheorem{lemma}[theorem]{Lemma}
\newtheorem{corollary}[theorem]{Corollary}

\theoremstyle{definition}
\newtheorem{definition}[theorem]{Definition}
\newtheorem{example}[theorem]{Example}

% ============================================================================
% ACTUS specification notation.
%
% These macros mirror the notation of the ACTUS Technical Specification
% (ACTUS Financial Research Foundation, v1.1, CC-BY-SA 4.0) so that the
% blueprint can reproduce the standard's schedule / state-transition / payoff
% formulas verbatim and link each to its Lean counterpart.  Names are kept
% close to the upstream source to ease cross-checking.
% ============================================================================

\newcommand{\Real}{\mathbb{R}}
\newcommand{\Nat}{\mathbb{N}}

\newcommand{\attr}[1]{\texttt{#1}}              % contract attribute (plain; math-safe)
% Linked contract attribute: \hyperlink to the term's \hypertarget in the Data
% Dictionary.  gen_glossary.py rewrites prose (text-mode) \attr{CODE} to
% \attrl{CODE} only where CODE has a dictionary entry; it is never emitted
% inside $...$, so it stays MathJax-safe in the web build.
\newcommand{\attrl}[1]{\hyperlink{term:#1}{\texttt{#1}}}
\newcommand{\svar}[2]{\ensuremath{\mathbf{#1}_{#2}}}  % state variable  Nt_{t}
% Plain/primed convention inside a state transition: a plain symbol is the
% pre-event value (the standard's t^-), a primed symbol the post-event value
% (the standard's t^+).  See the "Contract events, STF and POF" notes.
\newcommand{\sv}[1]{\ensuremath{\mathbf{#1}}}         % state var, pre-event value
\newcommand{\svn}[1]{\ensuremath{\mathbf{#1}'}}       % state var, post-event value
% The interest accrued on the outstanding notional since the last state date,
% Y(Sd,t)*Ipnr*Nt.  It recurs in nearly every transition; abbreviated here and
% defined once in the notation section.  \Delta is core LaTeX (plasTeX-safe).
\newcommand{\dipac}{\ensuremath{\Delta\mathbf{Ipac}}}
\newcommand{\nul}{\varnothing}                   % the null / undefined value (techspec $\undef$)
\newcommand{\sgn}{R(\attr{CNTRL})}               % contract-role sign

\newcommand{\stf}[2]{\ensuremath{\mathrm{STF}\_#1\_#2()}}  % named state-transition fn
\newcommand{\pof}[2]{\ensuremath{\mathrm{POF}\_#1\_#2()}}  % named payoff fn

% Fallback for plasTeX, which does not provide hyperref's \texorpdfstring;
% in the print build hyperref already defines it, so \providecommand is a no-op.
\providecommand{\texorpdfstring}[2]{#1}
\newcommand{\dfl}[1]{D(\mathbf{Prf}_{#1})}

\newcommand{\sdl}[3]{S(#1,#2,#3)}                % schedule
\newcommand{\sdll}[4]{S(#1,#2,#3,#4)}
\newcommand{\vsdl}[3]{\vec{S}(#1,#2,#3)}         % array schedule
\newcommand{\yfr}[2]{Y(#1,#2)}                   % year fraction
\newcommand{\yfrfunc}{Y}
\newcommand{\ann}[5]{A(#1,#2,#3,#4,#5)}          % annuity amount
\newcommand{\annfunc}{A}

\newcommand{\obs}[3]{O^{#1}(#2,#3)}              % risk-factor observer
\newcommand{\obsfull}[5]{O^{#1}(#2,#3,#4,#5)}
\newcommand{\obsfunc}[1]{O^{#1}}
\newcommand{\cldev}[3]{U^{ev}(#1,#2 \mid\{#3\})} % child-contract observer
\newcommand{\cldsv}[4]{U^{sv}(#1,#2,\svar{#3}{} \mid\{#4\})}
\newcommand{\cldsvs}[3]{U^{sv}(#1,#2,\svar{#3}{})}
\newcommand{\cldca}[2]{U^{ca}(#1,#2)}
\newcommand{\cldfunc}[1]{U^{#1}}

\newcommand{\tmax}{t^{max}}
\newcommand{\tev}[1]{\tau(#1)}                   % event time
\newcommand{\fev}[1]{f(#1)}                      % event type
\newcommand{\payoff}[2]{F(#1,#2)}                % canonical payoff

% The three normative tables that document each contract type --- the
% Contract Schedule, the State-Variable Initialization, and the State
% Transition / Payoff Functions --- are defined per output format:
% longtable for the print build (macros/print.tex, page-breakable) and
% plain tabular for the web build (macros/web.tex, rendered to HTML by
% plasTeX, which does not support longtable).  Both accept the same body:
% rows of three cells separated by & and terminated by \\ \hline.